% Draft supplemental material; all tables are editable inputs in sections/tables/.
\clearpage
\onecolumn
\WIPsection{Detailed Results: Draft Supplement}
\WIPsubsection{Occlusion and Procedure Summaries}
The tables retain the recording-level occlusion and timing summaries. Table~\ref{tab:suture-by-pig} preserves the numerical information from the removed stitch-overview figure as editable text.

% Editable table: update column specification, header, and row cells together.
\begin{table}[htbp]
\centering
\caption{Visual field obstruction.}
\label{tab:occlusion}
\small
\begin{tabularx}{\linewidth}{@{}Xrr@{}}
\toprule
Measure & 3 mm & 8 mm \\
\midrule
Recordings, $n$ & 3 & 4 \\
Frames in density plot, $n$ & 28,822 & 21,102 \\
Median of recording medians, \% & 10.39 & 23.93 \\
Range of recording medians, \% & 10.25--11.31 & 22.86--26.50 \\
Median of filtered recording medians, \% & 10.71 & 24.42 \\
Pooled filtered-frame median, \% & 11.0 & 25.1 \\
\bottomrule
\end{tabularx}
\par\smallskip
\begin{minipage}{\linewidth}\footnotesize
Dark frames were excluded; filtered summaries also excluded frames with less than 1\% visual field obstruction. Case-level summaries weighted recordings equally; pooled summaries weighted frames equally.
\end{minipage}
\end{table}

% Editable table: update column specification, header, and row cells together.
\begin{table}[htbp]
\centering
\caption{Timing summaries for the main recordings.}
\label{tab:sutures}
\small
\begin{tabularx}{\linewidth}{@{}Xrr@{}}
\toprule
Measure & 3 mm & 8 mm \\
\midrule
Recordings, $n$ & 3 & 4 \\
Annotated stitches, $n$ & 23 & 39 \\
Mean of per-pig stitch means, min & 7.1 & 4.1 \\
Mean of per-pig gap means, min & 5.4 & 2.0 \\
Median procedure duration, min & 115.91 & 75.46 \\
Median dissection duration, min & 45.57 & 22.62 \\
Median suturing span, min & 76.60 & 51.63 \\
\bottomrule
\end{tabularx}
\par\smallskip
\begin{minipage}{\linewidth}\footnotesize
Per-pig means weighted pigs equally. Spans included pauses. The P01 final assistant interval was excluded; that procedure ended at knot tightening. Stitches were clustered within pigs.
\end{minipage}
\end{table}

% Plain editable TeX. To add/remove a column, change the column specification,
% header, every data row, and the notes below. No automatic CSV import is used.
% Numerical check: data/sutures/by_pig_timing.csv, derived from stitches.csv.
\begin{table}[htbp]
\centering\small
\caption{Per-pig suture timing retained from the stitch-overview figure.}
\label{tab:suture-by-pig}
\begin{tabularx}{\linewidth}{@{}Xlrrrr@{}}
\toprule
Recording & Tool & Stitches & Mean stitch (min) & Mean gap (min) & Suturing span (min) \\
\midrule
May 29 & 3 mm & 4 & 10.96 & 10.92 & 76.60 \\
June 15 & 3 mm & 6 & 6.20 & 2.82 & 51.87 \\
June 16 & 3 mm & 13 & 4.11 & 2.42 & 82.89 \\
\textbf{Group mean} & 3 mm & \textbf{7.67} & \textbf{7.09} & \textbf{5.39} & \textbf{70.45} \\
\midrule
August 18 & 8 mm & 7 & 3.95 & 1.65 & 37.77 \\
August 19 & 8 mm & 13 & 2.98 & 1.17 & 53.04 \\
October 7 & 8 mm & 8 & 6.08 & 3.96 & 77.00 \\
October 8 & 8 mm & 11 & 3.58 & 1.04 & 50.21 \\
\textbf{Group mean} & 8 mm & \textbf{9.75} & \textbf{4.15} & \textbf{1.95} & \textbf{54.50} \\
\bottomrule
\end{tabularx}
\par\smallskip
\begin{minipage}{\linewidth}\footnotesize
Stitch duration runs from first needle contact to knot tightening. Gap runs from knot tightening to the next stitch's start and includes phase transitions. Span runs from first needle contact to the last retained cut or knot endpoint. The May 29 final assistant interval is excluded. Bold rows are arithmetic means of the pig-level values; the stitch count in these rows is therefore a mean, not a total. Each pig contributes equally. All recordings are from 2026.
\end{minipage}
\end{table}


\clearpage
\WIPsubsection{Per-Suture Phase Timing}
Each completed suture contributed a needle-passage interval, a knot-tying interval, and an assistant interval where retained. The figure shows all observations because the durations were skewed (Fig.~\ref{fig:suture-phases}). These pooled stitch-level medians differ from the means of per-pig means in the main timing summary.

% Editable table: update column specification, header, and row cells together.
\begin{table}[htbp]
\centering\small
\caption{Per-suture phase timing in seconds. Values are pooled medians (25th--75th percentiles).}
\label{tab:suture-phases}
\begin{tabularx}{\linewidth}{@{}XrrXX@{}}
\toprule
Phase & 3 mm, $n$ & 8 mm, $n$ & 3 mm time & 8 mm time \\
\midrule
Needle passage & 23 & 39 & 121.8 (56.0--254.8) & 71.3 (46.1--122.0) \\
Knot tying & 23 & 39 & 147.1 (113.8--204.2) & 123.4 (105.4--147.0) \\
Assistant & 22 & 39 & 34.0 (29.3--50.7) & 29.1 (24.2--42.2) \\
\bottomrule
\end{tabularx}
\par\smallskip\footnotesize
The 3 mm observations came from three pigs and the 8 mm observations from four. One May 29 assistant interval was excluded. Needle passage ran from first tissue contact to the last annotated needle-through event; knot tying ended at the last throw tightened. The assistant interval ended at cutting and was not part of the reported stitch duration. No independence or statistical significance is inferred from the number of stitches.
\end{table}


\clearpage
\WIPsubsection{Repair Regions and Stitch Components}
Timing was separated into needle passage, knot tying, and the assistant interval. Table~\ref{tab:suture-regions} reports pooled stitch medians within annotated repair regions. Region labels were not assumed to identify nasal or oral closure layers.

% Editable table: update column specification, header, and row cells together.
\begin{table}[htbp]
\centering\small
\caption{Stitch components by repair region. Values are medians (25th--75th percentiles), in seconds.}
\label{tab:suture-regions}
\begin{tabularx}{\linewidth}{@{}l l r r *{4}{>{\raggedright\arraybackslash}X}@{}}
\toprule
Region & Tool & Pigs & Stitches & Total stitch & Needle passage & Knot tying & Assistant \\
\midrule
Muscle & 3 mm & 1 & 2 & 348.3 (323.0--373.6) & 208.2 (184.4--232.0) & 140.1 (138.7--141.6) & 52.8 (44.8--60.8) \\
Muscle & 8 mm & 3 & 5 & 338.5 (202.8--354.3) & 185.1 (94.3--207.9) & 130.5 (108.6--169.2) & 38.7 (28.3--52.3) \\
Soft palate & 3 mm & 3 & 13 & 335.0 (226.2--430.6) & 121.8 (56.5--253.9) & 178.4 (143.2--247.9) & 39.5 (32.3--65.6) \\
Soft palate & 8 mm & 4 & 18 & 211.7 (176.6--297.0) & 84.0 (46.8--132.7) & 129.8 (115.8--163.3) & 28.9 (24.0--42.5) \\
Hard palate & 3 mm & 2 & 8 & 206.1 (161.4--356.7) & 76.7 (54.2--167.2) & 113.8 (92.9--158.8) & 27.0 (19.4--32.5) \\
Hard palate & 8 mm & 4 & 16 & 173.9 (158.4--201.9) & 61.7 (40.5--71.6) & 111.4 (99.2--142.0) & 29.1 (23.8--40.4) \\
\bottomrule
\end{tabularx}
\par\smallskip\footnotesize
The soft-palate 3 mm assistant summary contains 12 observations, versus 13 stitches, after the May 29 exclusion. Other component counts equal the stitch counts. Stitch duration ends at knot tightening and excludes the assistant interval. Pooled medians weight pigs with more stitches more heavily; component medians need not sum to the total median.
\end{table}


\WIPsubsection{Knot Throws}
Throws were counted from knot video clips. Figure~\ref{fig:knot-throws-mode} shows the throw counts and the mode version of the timing distribution. Both groups had a median of three throws per knot. The pooled median of the knot-tying interval per counted throw was 52.6 seconds with 3 mm instruments and 43.5 seconds with 8 mm instruments (Table~\ref{tab:knot-throws}).

% Editable table: update column specification, header, and row cells together.
\begin{table}[htbp]
\centering\small
\caption{Knot-throw counts and timing. Values are pooled medians (25th--75th percentiles).}
\label{tab:knot-throws}
\begin{tabularx}{\linewidth}{@{}lrrXXX@{}}
\toprule
Tool & Pigs & Knots & Throws per knot & Knot interval (s) & Interval per throw (s) \\
\midrule
3 mm & 3 & 23 & 3.0 (2.0--3.0) & 147.1 (113.8--204.2) & 52.6 (37.9--84.2) \\
8 mm & 4 & 39 & 3.0 (3.0--3.0) & 123.4 (105.4--147.0) & 43.5 (35.8--51.6) \\
\bottomrule
\end{tabularx}
\end{table}


\begin{itemize}
\item The knot interval began at needle passage and included pulling suture through tissue. Dividing it by the throw count does not measure isolated throw execution time.
\item First-of-block knots had longer median intervals per throw: 65 versus 48 seconds for later knots with 3 mm instruments, and 44 versus 41 seconds with 8 mm instruments. Thread length was not recorded.
\end{itemize}

\clearpage
\WIPsubsection{Phase-Level Instrument Motion}
Motion analysis included 19 stitches from two 3 mm pigs and 39 stitches from four 8 mm pigs. May 29 had no rosbag. Table~\ref{tab:phase-motion} separates needle passage from knot tying and reports both phase totals and duration-normalised values.

% Editable table: update column specification, header, and row cells together.
\begin{table}[htbp]
\centering\small
\caption{Phase-level instrument motion. Values are pooled medians (25th--75th percentiles).}
\label{tab:phase-motion}
\begin{tabularx}{\linewidth}{@{}>{\raggedright\arraybackslash}p{.28\linewidth}*{4}{>{\raggedright\arraybackslash}X}@{}}
\toprule
Measure & Needle, 3 mm & Needle, 8 mm & Knot, 3 mm & Knot, 8 mm \\
\midrule
Shaft-end path (mm) & 394.9 (264.8--1039.4) & 427.0 (269.3--720.6) & 1204.4 (1033.7--1556.6) & 1072.0 (905.7--1262.4) \\
Roll (rad) & 20.5 (11.1--40.5) & 21.2 (14.7--35.1) & 35.7 (29.5--42.0) & 39.6 (33.7--46.9) \\
Wrist pitch + yaw (rad) & 28.7 (17.8--75.7) & 30.0 (17.1--41.0) & 57.7 (47.4--73.7) & 54.1 (43.8--63.1) \\
Jaw closures & 10.0 (6.5--21.5) & 14.0 (9.0--23.5) & 21.0 (16.0--25.5) & 23.0 (21.0--30.5) \\
Both instruments still (s) & 16.4 (6.8--29.2) & 8.9 (4.9--17.1) & 11.6 (6.2--14.8) & 9.2 (7.1--13.3) \\
Shaft-end path (mm/min) & 302.1 (283.8--340.6) & 359.3 (323.9--399.5) & 528.3 (468.8--649.6) & 522.9 (469.3--577.4) \\
Roll (rad/min) & 12.5 (10.6--14.9) & 19.8 (16.8--22.1) & 15.4 (13.9--16.9) & 19.6 (15.7--20.6) \\
Wrist pitch + yaw (rad/min) & 20.8 (18.5--23.9) & 22.9 (20.5--26.3) & 24.6 (23.1--26.9) & 25.4 (21.5--31.7) \\
Jaw closures per minute & 7.4 (6.4--8.5) & 11.8 (9.4--13.8) & 9.3 (8.0--10.0) & 11.8 (10.4--13.2) \\
Both instruments still (\%) & 14.3 (12.4--16.0) & 13.5 (10.2--15.8) & 7.1 (5.5--8.7) & 8.0 (6.1--9.9) \\
Whole-phase workspace (cm$^3$) & 1.5 (1.0--2.4) & 2.1 (1.0--3.3) & 6.9 (3.6--10.4) & 5.9 (4.0--7.7) \\
First-60-s workspace (cm$^3$) & 1.0 (0.6--1.3) [n=12] & 1.6 (1.1--2.2) [n=24] & 3.7 (2.1--5.8) [n=19] & 3.1 (2.2--3.9) [n=38] \\
Spread about own centre (mm) & 4.4 (4.1--4.9) & 5.1 (4.6--6.0) & 6.6 (5.5--7.4) & 6.6 (6.1--7.1) \\
\bottomrule
\end{tabularx}
\par\smallskip\footnotesize
Both instruments were summed for path length, rotations, and jaw closures. Workspace is the sum of separate shaft-end convex-hull volumes in instrument-relative frames, not a registered anatomical workspace or tool-tip volume. Spread is averaged across instruments. First-60-s volumes exclude phases shorter than 60 seconds; their counts are shown in brackets. All other columns contain 19 or 39 observations, respectively.
\end{table}


\begin{itemize}
\item Signals were averaged into 0.05-second bins. Jaw closure meant closing below 0.1 rad after opening above 0.3 rad; this does not establish a successful tissue grasp.
\item Both instruments were classified as still when each shaft end moved below 2 mm/s and wrist rotation below 0.1 rad/s.
\item \textbf{[Confirm synchronization and kinematic assumptions before interpreting these motion estimates as efficiency.]}
\item The analysis repository also contains preliminary insertion-effort estimates without a force sensor. These have not been presented as validated tissue force.
\end{itemize}
